\bigskip\Needspace{10\baselineskip}
\pdfbookmark[1]{Fortgesetzte Nachrichten. Ceres Ferdinandea. April 1802}{ceres180204}
\begin{center}
{\small FORTGESETZTE NACHRICHTEN ÜBER DEN NEUEN HAUPTPLANETEN.\par}
CERES FERDINANDEA.\par
\medskip\hrule\medskip
Monatliche Correspondenz Band V. S.\ 379—404. April 1802.\par
\medskip\hrule
\end{center}\medskip
—\quad —\quad —\quad Die VII. elliptische Bahn des \textit{neuen Planeten}, die \textsc{Gauß} bestimmt hat, [S.\ 207 dieses Bandes] stimmt weiterhin genau mit unseren fortgesetzten Beobachtungen dieses planetarischen Himmelskörpers überein, wie die folgende Berechnung und der Vergleich von Dr.\ \textsc{Gauß} hinreichend zeigen.
\begin{center}
{\small\setlength{\tabcolsep}{5pt}
\begin{tabular}{r r|r|r|r|r}
\multicolumn{2}{c|}{\shortstack{Seeberg\\1802}}&\shortstack{Berechnete\\Rektaszension\\$\CeresSymbol$}&\shortstack{Unter-\\schied}&\shortstack{Berechnete\\nördliche\\Deklination}&\shortstack{Unter-\\schied}\\\hline
Februar&19&$187^0\ 58'\ 23''8$&$-4''1$&$14^0\ 20'\ 24''0$&$+21''1$\\
—&26&$187\quad7\ 34.4$&$+1.0$&$15\quad9\ 20.1$&$+25.8$\\
—&27&$186\ 58\ 53.2$&$+5.8$&$15\ 16\ 20.2$&$+26.3$\\
—&28&$186\ 49\ 51.9$&$+3.7$&$15\ 23\ 19.1$&$+18.1$\\
März&1&$186\ 40\ 31.4$&$+3.5$&$15\ 30\ 16.0$&$+34.9$\\
—&2&$186\ 30\ 52.2$&$+3.3$&$15\ 37\ 10.6$&$\ .\ .\ .$\\
—&3&$186\ 20\ 55.1$&$+2.6$&$15\ 44\quad2.0$&$+19.7$
\end{tabular}}
\end{center}
Die Rektaszensionen stimmen also weiterhin ausgezeichnet überein, und der Fehler in der Deklination scheint sich wenig oder gar nicht geändert zu haben. Lässt man die Beobachtung vom 1. März weg, so läge er im Mittel aus fünf Beobachtungen bei $+22''$. —\quad —\quad —

Oberamtmann Dr.\ \textsc{Schröter} beobachtet den Durchmesser der \textit{Ceres} weiterhin regelmäßig. Hier folgen einige seiner Beobachtungen, die Dr.\ \textsc{Gauß} berechnet und uns freundlicherweise mitgeteilt hat.
\begin{center}
{\small\setlength{\tabcolsep}{4pt}
\begin{tabular}{r l r l r}
\multicolumn{3}{c}{Scheinbarer Durchmesser}&\multicolumn{2}{c}{\shortstack{Berechneter Durchmesser\\im Abstand 1.}}\\[5pt]
1802 Januar 10&$\ .\ .\ .\ .$&$2''500$&$\ .\ .\ .\ .\ .\ .$&$5''16$\\
25&$\ .\ .\ .\ .$&$2.514$&$\ .\ .\ .\ .\ \raisebox{.5ex}{.}\ .$&$4.77$\\
26&$\ .\ .\ .\ .$&$2.687$&$\ .\ .\ .\ .\ .\ .$&$5.07$\\
28&$\ .\ .\ .\ .$&$2.774$&$\ .\ .\ .\ .\ .\ .$&$5.18$\\
31&$\ .\ .\ .\ .$&$2.930$&$\ .\ .\ .\ .\ .\ .$&$5.38$\\
Februar 5&$\ .\ .\ .\ .$&$3.120$&$\ .\ .\ .\ .\ .\ .$&$5.59$\\
10&$\ .\ .\ .\ .$&$3.342$&$\ .\ .\ .\ .\ .\ .$&$5.90$
\end{tabular}}
\end{center}
\SourceLeaf{215}%

Das Mittel wäre $5''29=0.308$ der Erde oder 529 geografische Meilen, einschließlich des. atmosphärischen Nebels. Nach \textsc{Schröter}s Messung der scharf begrenzten Scheibe betrüge deren Durchmesser nur $3''44$ in der Entfernung $=1$, also fast genau $\tfrac15$ des Erddurchmessers. Damit wäre sie beträchtlich kleiner als unser Mond. Kein Wunder also, dass uns ein solches Planetchen so lange verborgen geblieben ist!

Da die VII. Ellipse für die \textit{Ceres}-Bahn von \textsc{Gauß} die Beobachtungen dieses Planeten weiterhin recht genau wiedergibt, teilen wir unseren astronomischen Lesern noch rechtzeitig eine Fortsetzung der \textit{Ephemeride} ihres Laufs für die nächsten drei Monate mit. Dr.\ \textsc{Gauß} hat die Berechnung freundlicherweise übernommen. Freilich ist kaum zu glauben, dass \textit{Ceres} bis dahin sichtbar bleibt. Schon jetzt, am 27. und 28. März, beobachten wir bei diesem Planeten eine so rasche Veränderung des Lichts, dass wir geneigt sind, sie ganz anderen Ursachen zuzuschreiben als dem Zustand unserer Atmosphäre oder der Entfernung des Planeten von der Erde. Astronomen, die mit Äquatorial- und parallaktischen Instrumenten ausgestattet sind und nichts unversucht lassen wollen, können sich damit die immer schwieriger werdende Suche nach Ceres erheblich erleichtern.
\begin{center}
Stellungen der \textit{Ceres} im April, Mai und Juni 1802 für Mitternacht, mittlere Seeberger Zeit.\par\medskip
{\fontsize{9}{11}\selectfont\setlength{\tabcolsep}{2pt}
\begin{tabular}{r r|r|r|r@{\qquad}r r|r|r|r}
\multicolumn{2}{c|}{}&\shortstack{R.A. $\CeresSymbol$\\in Grad}&\shortstack{Nördliche\\Dekl.}&\shortstack{R.A. $\CeresSymbol$\\in Zeit}&\multicolumn{2}{c|}{}&\shortstack{R.A. $\CeresSymbol$\\in Grad}&\shortstack{Nördliche\\Dekl.}&\shortstack{R.A. $\CeresSymbol$\\in Zeit}\\\hline
April&21&$177^0\ 14'$&$17^0\ 52'$&$11^{\mathrm u}\ 48^{\mathrm m}\ 55^{\mathrm s}$&Mai&27&$177^0\ 14'$&$14^0\ 13'$&$11^{\mathrm u}\ 48^{\mathrm m}\ 57^{\mathrm s}$\\
&24&$176\ 56$&$17\ 42$&$11\ 47\ 43$&&30&$177\ 34$&$13\ 47$&$11\ 50\ 16$\\
&27&$176\ 41$&$17\ 30$&$11\ 46\ 44$&Juni&2&$177\ 57$&$13\ 21$&$11\ 51\ 47$\\
&30&$176\ 29$&$17\ 17$&$11\ 45\ 58$&&5&$178\ 22$&$12\ 53$&$11\ 53\ 27$\\
Mai&3&$176\ 21$&$17\quad2$&$11\ 45\ 25$&&8&$178\ 49$&$12\ 25$&$11\ 55\ 17$\\
&6&$176\ 17$&$16\ 45$&$11\ 45\quad6$&&11&$179\ 19$&$11\ 56$&$11\ 57\ 16$\\
&9&$176\ 15$&$16\ 27$&$11\ 45\quad2$&&14&$179\ 51$&$11\ 27$&$11\ 59\ 24$\\
&12&$176\ 17$&$16\quad8$&$11\ 45\ 10$&&17&$180\ 25$&$10\ 57$&$12\quad1\ 40$\\
&15&$176\ 23$&$15\ 47$&$11\ 45\ 27$&&20&$181\quad1$&$10\ 27$&$12\quad4\quad4$\\
&18&$176\ 31$&$15\ 25$&$11\ 46\quad5$&&23&$181\ 39$&$9\ 56$&$12\quad6\ 35$\\
&21&$176\ 43$&$15\quad2$&$11\ 46\ 50$&&26&$182\ 18$&$9\ 25$&$12\quad9\ 13$\\
&24&$176\ 57$&$14\ 38$&$11\ 47\ 48$&&29&$183\quad0$&$8\ 54$&$12\ 11\ 59$
\end{tabular}}
\end{center}

\textit{Ceres} kommt zum Stillstand und wird am 2. Mai in der Länge bei $5^{\mathrm z}\ 19^0\ 45'$ und am 9. Mai in der R.A. bei $176^0\ 15'$ wieder rechtläufig. Die rückläufige Bewegung beträgt diesmal in der Länge $13^0\ 10'$ und dauert 92 Tage. In der R.A. beträgt sie $12^0\ 27'$ und dauert 93 Tage. Das Verhältnis der Lichtstärke des Planeten hat Dr.\ \textsc{Gauß} unter Vernachlässigung der Größe der Phase wie folgt berechnet.
\begin{center}
{\small\setlength{\tabcolsep}{5pt}
\begin{tabular}{r l r|r|r|r}
\multicolumn{3}{c|}{}&\multicolumn{2}{c|}{Abstand von der}&\shortstack{Licht-\\stärke}\\
\multicolumn{3}{c|}{}&$\odot$&$\EarthSymbol$&\\\hline
1801&Dezember&7&$2.5453$&$2.4892$&$0.4227$\\
1802&März&16&$2.5706$&$1.6024$&$1.0000$\\
&Mai&6&$2.5961$&$1.9039$&$0.6945$\\
&Juni&29&$2.6305$&$2.5735$&$0.3703$
\end{tabular}}
\end{center}
\begin{flushright}[\textsc{von Zach}.]\end{flushright}
